\documentclass[a4paper, 12pt]{report}

%% Declare the packages used in the template
\usepackage{amssymb,amsmath,amsfonts,eurosym,geometry,ulem,graphicx,caption,color,setspace,sectsty,comment,footmisc,caption,pdflscape,subfigure,array,authblk,lipsum,float,graphicx,blindtext}
\graphicspath{{attachment/}} % Path to the folder containing the image files
\usepackage[hidelinks]{hyperref} % Make cross-references and hyperlinks clickable without boders around

\usepackage{subfiles} % Best loaded last in the preamble

\usepackage{natbib}         % For creating bibliography.

\usepackage{pythonhighlight} %For typing Python code
\usepackage[Sonny]{fncychap} %For Chapter style
    \ChNameUpperCase
    \ChTitleVar{\Large\sc}

\usepackage{listings} %For adding other language code
\usepackage{upquote}
\usepackage{xcolor}

%%%% Set up styles and colour for programming code
\definecolor{codegreen}{rgb}{0,0.6,0}
\definecolor{codegray}{rgb}{0.5,0.5,0.5}
\definecolor{codepurple}{rgb}{0.58,0,0.82}
\definecolor{backcolour}{rgb}{0.95,0.95,0.92}

\lstdefinestyle{mystyle}{
    backgroundcolor=\color{backcolour},   
    commentstyle=\color{codegreen},
    keywordstyle=\color{magenta},
    numberstyle=\tiny\color{codegray},
    stringstyle=\color{codepurple},
    basicstyle=\ttfamily\footnotesize,
    breakatwhitespace=false,         
    breaklines=true,                 
    captionpos=b,                    
    keepspaces=true,                 
    numbers=left,                    
    numbersep=5pt,                  
    showspaces=false,                
    showstringspaces=false,
    showtabs=false,                  
    tabsize=2
}

\lstset{style=mystyle}
    
\normalem

\onehalfspacing
\newtheorem{theorem}{Theorem}
\newtheorem{corollary}[theorem]{Corollary}
\newtheorem{lemma}[theorem]{Lemma}
\newtheorem{proposition}{Proposition}
\newenvironment{proof}[1][Proof]{\noindent\textbf{#1.} }{\ \rule{0.5em}{0.5em}}

\newtheorem{hyp}{Hypothesis} 
\newtheorem{subhyp}{Hypothesis}

\newcommand{\red}[1]{{\color{red} #1}}
\newcommand{\blue}[1]{{\color{blue} #1}}

\newcolumntype{L}[1]{>{\raggedright\let\newline\\arraybackslash\hspace{0pt}}m{#1}}
\newcolumntype{C}[1]{>{\centering\let\newline\\arraybackslash\hspace{0pt}}m{#1}}
\newcolumntype{R}[1]{>{\raggedleft\let\newline\\arraybackslash\hspace{0pt}}m{#1}}

\geometry{left=1.0in,right=1.0in,top=1.0in,bottom=1.0in}
%%%%%%%%%%%%%%%%%%%%%%%%%%%%%%%%%%%%%%%%%%%%%%%%%%%%%%%%%%%%%%%%%

\begin{document}

%%%% COVER PAGE (include Abstract) %%%%%
\begin{titlepage}
    \begin{center}
        \vspace*{1cm}
            \text{\Large MPP THESIS (\the\year{})}
            
        \vspace{4cm}
            \textbf{\LARGE Thesis Title}
            
        \vspace{1 cm}
            \includegraphics[width=0.4\textwidth]{appp_sym.png}

        \vspace{2.5cm}
           \today
            
        \vspace{.5cm}    
            \text{\Large Author Name}
            
        \vspace{1.5cm}
            \Large
                Asian Public Policy Program\\
                School of International and Public Policy\\
                Hitotsubashi University
   
    \setcounter{page}{0}
    \thispagestyle{empty}            
    \end{center}
\end{titlepage}


\pagebreak \newpage
%%%%%%%%%%%%%%%%%%%%%%%%%%%%%%%%%%%%%%%%%%%%%%%%%%%%%%%%%
\clearpage %Empty page

%%%%%%%%%%%%%%%% GENERAL INFORMATION %%%%%%%%%%%%%%%%%%%%%%%%%
% You can modify sections and subsection to fit with your work.

%%%% ABSTRACT %%%%
    \doublespacing
    
\chapter*{Abstract} \label{sec:abstract}
\pagenumbering{roman}
\lipsum[1-2]


%%%%%%%%%%%%%%%%%%%%%%%%%
\pagebreak \newpage

%%%% ACKNOWLEDGMENTS %%%%
    \doublespacing
\chapter*{Acknowledgements} \label{sec:acknowledgment}


%%%%%%%%%%%%%%%%%%%%%%%%%
\pagebreak \newpage

%%%% TABLE OF CONTENTS %%%%
    \doublespacing
\tableofcontents

%%%%%%%%%%%%%%%%%%%%%%%%%
\pagebreak \newpage

%%%% LIST OF FIGURES %%%%
\chapter*{List of Figures} \label{sec:figures}
    \addcontentsline{toc}{section}{Figures}

%%%% LIST OF TABLES %%%%
    \vspace{1.5cm}
\chapter*{List of Tables} \label{sec:tables}
    \addcontentsline{toc}{section}{Tables}

%%%% LIST OF APPENDICES %%%%
    \vspace{1.5cm}
    \addcontentsline{toc}{section}{Appendix A}

%%%%%%%%%%%%%%%%%%%%%%%%%
\pagebreak \newpage

%%%% LIST OF ABBREVIATIONS %%%%
\chapter*{List of Abbreviations} \label{sec:abbres}


%%%%%%%%%%%%%%%%%%%%%%%%%
\pagebreak \newpage

%%%% GLOSSARY %%%%
\chapter*{Glossary} \label{sec:glossary}

%%%%%%%%%%%%%%%%%%%%%%%%%%%%%%%%%%%%%%%%%%%%%%%%%%%%%%%%%

%%%%%%%%%%%%%%%% THESIS CONTENT %%%%%%%%%%%%%%%%%%%%%%%%%
% You can modify sections and subsection to fit with your work.
\pagebreak \newpage
\pagenumbering{arabic}

%%%% INTRODUCTION %%%%
\chapter{Introduction} \label{sec:introduction}

    \section*{Statement of the problem}


    
    \section*{Objectives}
    
    \section*{Organization of the thesis}



%%%%%%%%%%%%%%%%%%%%%%%%%
\pagebreak \newpage

%%%% BACKGROUND INFORMATION %%%%
\chapter{Background Information} \label{sec:background}



%%%%%%%%%%%%%%%%%%%%%%%%%
\pagebreak \newpage

%%%% LITERATURE REVIEW %%%%
\chapter{Literature Review} \label{sec:literature}
    Citation example \\
     For articles: \citep{Blanchard19} \\
     For books::  \cite[see][chap 2]{latexcompanion}

    \begin{hyp}
        The better the answer, the higher the score.
    \end{hyp}
    
    \stepcounter{hyp}
    \setcounter{subhyp}{0}
    \begin{subhyp}
        A higher score is positively correlated with shoe-size.
    \end{subhyp}
    
    \begin{subhyp}
        A lower score is negatively correlated with shoe-size.
    \end{subhyp}
    
    \begin{hyp}
        Pythagoras had something to say.
    \end{hyp}
    
    \setcounter{subhyp}{0} % Don't execute                      \stepcounter{hyp} in this case.
    \begin{subhyp}
        Pythagoras contributed something to geometry.
    \end{subhyp}
    
    \begin{subhyp}
        \emph{Pythagoras} wasn't a single person but a group of like-minded philosophers and mathematicians.
    \end{subhyp}
    
%%%%%%%%%%%%%%%%%%%%%%%%%
\pagebreak \newpage

%%%% THEORETICAL FRAMEWORK %%%%
\chapter{Theoretical Framework} \label{sec:theory}

        Examples for insert Theorem, Corollary, and Lemma
        \begin{itemize}
            \item Theorem
                \begin{theorem}
                    Let \(f\) be a function whose derivative exists in every point, then \(f\) is a continuous function.
                \end{theorem}

                \begin{theorem}[Pythagorean theorem]
                    \label{pythagorean}
                    This is a theorem about right triangles and can be summarised in the next equation 
                    \[ x^2 + y^2 = z^2 \]
                \end{theorem}

            And a consequence of theorem \ref{pythagorean} is the statement in the next corollary.

            \item Corollary
                \begin{corollary}
                    There's no right rectangle whose sides measure 3cm, 4cm, and 6cm.
                \end{corollary}

            \item Lemma
            
                \begin{lemma}
                    Given two line segments whose lengths are \(a\) and \(b\) respectively there is a real number \(r\) such that \(b=ra\).
                \end{lemma}
        \end{itemize}    

%%%%%%%%%%%%%%%%%%%%%%%%%
\pagebreak \newpage

%%%% EMPIRICAL FRAMEWORK %%%%
\chapter{Empirical Framework} \label{sec:model}

    Figure insertation example
        \begin{figure}[H]
            \centering
            \includegraphics[width=80mm]{attachment/appp_sym.png}
            \caption{APPP Symbol}
            \label{fig1}
            \end{figure}  

%%%%%%%%%%%%%%%%%%%%%%%%%
\pagebreak \newpage

%%%% EMPIRICAL RESULTS %%%%
\chapter{Empirical Results} \label{sec:result}


%%%%%%%%%%%%%%%%%%%%%%%%%
\pagebreak \newpage

%%%% CONCLUSIONS AND POLICY IMPLICATIONS %%%%
\chapter{Conclusion} \label{sec:conclusion}

    \section{Conclusions}
    
    \section{Policy Implications}
    
    \section{Limitations and Further Research}

%%%%%%%%%%%%%%%%%%%%%%%%%
\pagebreak \newpage

%%%% APPENDICES %%%%
\chapter*{Appendices} \label{sec:appendix}

\appendix

\section*{Appendix A: Stata code}\label{appendixa}
    
\begin{lstlisting}
* Run a sharp RDD with a second-order polynomial term rdrobust y x, c(0) p(2)
* Run a fuzzy RDD
* We don't have a fuzzy RDD in this data, but let's create one, where
* probability of treatment jumps from 20% to 60% at the cutoff
    g treatment = (runiform() < .2)*(x < 0) + (runiform() < .6)*(x >= 0)
    rdrobust y x, c(0) fuzzy(treatment)
* Generate a standard RDD plot with a polynomial of 2 (default is 4)
    rdplot y x, c(0) p(2)
\end{lstlisting}
    
\section*{Appendix B. R Programming Code}\label{appendixb}
\subsection*{R Programming Code}

\begin{lstlisting}[language=R]
fib <- function(n) {
    if (n < 2)
        n
    else
        fib(n - 1) + fib(n - 2)
    }
fib(10) # => 55
\end{lstlisting}

\section*{Appendix C: Python code}\label{appendixc}
\begin{python}
##################### Start of code ###########################
# Problem 2.1. Game Theory with Nashpy package for 2 players and 2 strategies
import nashpy as nash
import numpy as np
import fractions
        
# NumPy: Convert decimals to fractions
np.set_printoptions(formatter={'all':lambda x: str(fractions.Fraction(x).limit_denominator())})
        
# Create the payoff matrix
R = np.array([[-2, 3],
             [-1, 0]]) # R is the Row player (Passenger)
        
C = np.array([[-1, -9],
             [-2,  0]]) # C is the Col player (Inspector)
        
game21 = nash.Game(R,C)
game21
        
# Find the Nash Equilibrium with Support Enumeration
equilibria = game21.support_enumeration()
list(equilibria)
##################### End of code ###########################

# Result
[(array([1/5, 4/5]), array([3/4, 1/4]))]
\end{python} 


%%%%%%%%%%%%%%%%%%%%%%%%%
\pagebreak \newpage

%%%% REFERENCES %%%%
% This part is fixed setup - DON'T CHANGE
    \singlespacing
    \bibliographystyle{ecta}
    \bibliography{Bibs} %Imports the file sample.bib that contains bibliography sources
    
    %\printbibliography %Prints bibliography
%%%%% THE END OF THESIS CONTENT %%%%%%%%%%%%%%%%%%%%%%

\end{document}